% Options for packages loaded elsewhere
\PassOptionsToPackage{unicode}{hyperref}
\PassOptionsToPackage{hyphens}{url}
\documentclass[
]{article}
\usepackage{xcolor}
\usepackage[margin=1in]{geometry}
\usepackage{amsmath,amssymb}
\setcounter{secnumdepth}{-\maxdimen} % remove section numbering
\usepackage{iftex}
\ifPDFTeX
  \usepackage[T1]{fontenc}
  \usepackage[utf8]{inputenc}
  \usepackage{textcomp} % provide euro and other symbols
\else % if luatex or xetex
  \usepackage{unicode-math} % this also loads fontspec
  \defaultfontfeatures{Scale=MatchLowercase}
  \defaultfontfeatures[\rmfamily]{Ligatures=TeX,Scale=1}
\fi
\usepackage{lmodern}
\ifPDFTeX\else
  % xetex/luatex font selection
\fi
% Use upquote if available, for straight quotes in verbatim environments
\IfFileExists{upquote.sty}{\usepackage{upquote}}{}
\IfFileExists{microtype.sty}{% use microtype if available
  \usepackage[]{microtype}
  \UseMicrotypeSet[protrusion]{basicmath} % disable protrusion for tt fonts
}{}
\makeatletter
\@ifundefined{KOMAClassName}{% if non-KOMA class
  \IfFileExists{parskip.sty}{%
    \usepackage{parskip}
  }{% else
    \setlength{\parindent}{0pt}
    \setlength{\parskip}{6pt plus 2pt minus 1pt}}
}{% if KOMA class
  \KOMAoptions{parskip=half}}
\makeatother
\usepackage{color}
\usepackage{fancyvrb}
\newcommand{\VerbBar}{|}
\newcommand{\VERB}{\Verb[commandchars=\\\{\}]}
\DefineVerbatimEnvironment{Highlighting}{Verbatim}{commandchars=\\\{\}}
% Add ',fontsize=\small' for more characters per line
\usepackage{framed}
\definecolor{shadecolor}{RGB}{248,248,248}
\newenvironment{Shaded}{\begin{snugshade}}{\end{snugshade}}
\newcommand{\AlertTok}[1]{\textcolor[rgb]{0.94,0.16,0.16}{#1}}
\newcommand{\AnnotationTok}[1]{\textcolor[rgb]{0.56,0.35,0.01}{\textbf{\textit{#1}}}}
\newcommand{\AttributeTok}[1]{\textcolor[rgb]{0.13,0.29,0.53}{#1}}
\newcommand{\BaseNTok}[1]{\textcolor[rgb]{0.00,0.00,0.81}{#1}}
\newcommand{\BuiltInTok}[1]{#1}
\newcommand{\CharTok}[1]{\textcolor[rgb]{0.31,0.60,0.02}{#1}}
\newcommand{\CommentTok}[1]{\textcolor[rgb]{0.56,0.35,0.01}{\textit{#1}}}
\newcommand{\CommentVarTok}[1]{\textcolor[rgb]{0.56,0.35,0.01}{\textbf{\textit{#1}}}}
\newcommand{\ConstantTok}[1]{\textcolor[rgb]{0.56,0.35,0.01}{#1}}
\newcommand{\ControlFlowTok}[1]{\textcolor[rgb]{0.13,0.29,0.53}{\textbf{#1}}}
\newcommand{\DataTypeTok}[1]{\textcolor[rgb]{0.13,0.29,0.53}{#1}}
\newcommand{\DecValTok}[1]{\textcolor[rgb]{0.00,0.00,0.81}{#1}}
\newcommand{\DocumentationTok}[1]{\textcolor[rgb]{0.56,0.35,0.01}{\textbf{\textit{#1}}}}
\newcommand{\ErrorTok}[1]{\textcolor[rgb]{0.64,0.00,0.00}{\textbf{#1}}}
\newcommand{\ExtensionTok}[1]{#1}
\newcommand{\FloatTok}[1]{\textcolor[rgb]{0.00,0.00,0.81}{#1}}
\newcommand{\FunctionTok}[1]{\textcolor[rgb]{0.13,0.29,0.53}{\textbf{#1}}}
\newcommand{\ImportTok}[1]{#1}
\newcommand{\InformationTok}[1]{\textcolor[rgb]{0.56,0.35,0.01}{\textbf{\textit{#1}}}}
\newcommand{\KeywordTok}[1]{\textcolor[rgb]{0.13,0.29,0.53}{\textbf{#1}}}
\newcommand{\NormalTok}[1]{#1}
\newcommand{\OperatorTok}[1]{\textcolor[rgb]{0.81,0.36,0.00}{\textbf{#1}}}
\newcommand{\OtherTok}[1]{\textcolor[rgb]{0.56,0.35,0.01}{#1}}
\newcommand{\PreprocessorTok}[1]{\textcolor[rgb]{0.56,0.35,0.01}{\textit{#1}}}
\newcommand{\RegionMarkerTok}[1]{#1}
\newcommand{\SpecialCharTok}[1]{\textcolor[rgb]{0.81,0.36,0.00}{\textbf{#1}}}
\newcommand{\SpecialStringTok}[1]{\textcolor[rgb]{0.31,0.60,0.02}{#1}}
\newcommand{\StringTok}[1]{\textcolor[rgb]{0.31,0.60,0.02}{#1}}
\newcommand{\VariableTok}[1]{\textcolor[rgb]{0.00,0.00,0.00}{#1}}
\newcommand{\VerbatimStringTok}[1]{\textcolor[rgb]{0.31,0.60,0.02}{#1}}
\newcommand{\WarningTok}[1]{\textcolor[rgb]{0.56,0.35,0.01}{\textbf{\textit{#1}}}}
\usepackage{graphicx}
\makeatletter
\newsavebox\pandoc@box
\newcommand*\pandocbounded[1]{% scales image to fit in text height/width
  \sbox\pandoc@box{#1}%
  \Gscale@div\@tempa{\textheight}{\dimexpr\ht\pandoc@box+\dp\pandoc@box\relax}%
  \Gscale@div\@tempb{\linewidth}{\wd\pandoc@box}%
  \ifdim\@tempb\p@<\@tempa\p@\let\@tempa\@tempb\fi% select the smaller of both
  \ifdim\@tempa\p@<\p@\scalebox{\@tempa}{\usebox\pandoc@box}%
  \else\usebox{\pandoc@box}%
  \fi%
}
% Set default figure placement to htbp
\def\fps@figure{htbp}
\makeatother
\setlength{\emergencystretch}{3em} % prevent overfull lines
\providecommand{\tightlist}{%
  \setlength{\itemsep}{0pt}\setlength{\parskip}{0pt}}
\usepackage{bookmark}
\IfFileExists{xurl.sty}{\usepackage{xurl}}{} % add URL line breaks if available
\urlstyle{same}
\hypersetup{
  hidelinks,
  pdfcreator={LaTeX via pandoc}}

\author{}
\date{\vspace{-2.5em}}

\begin{document}

\begin{Shaded}
\begin{Highlighting}[]
\FunctionTok{getwd}\NormalTok{()}
\end{Highlighting}
\end{Shaded}

\begin{verbatim}
## [1] "/Users/johann/Documents/labstuff/dataviz/pydataviz/dataviz"
\end{verbatim}

\begin{Shaded}
\begin{Highlighting}[]
\NormalTok{path }\OtherTok{\textless{}{-}} \StringTok{"assets/img/RStudio.png"}
\FunctionTok{file.exists}\NormalTok{(path)}
\end{Highlighting}
\end{Shaded}

\begin{verbatim}
## [1] TRUE
\end{verbatim}

\section*{(PART) Get}\label{part-get}
\addcontentsline{toc}{section}{(PART) Get}

\section{R basics}\label{r-basics}

This chapter provides a quick introduction to the programming language R
and to using the software RStudio.

\subsection{Rstudio}\label{rstudio}

Rstudio is a software that allows to program in R and interactively
analyse data with R. It succinctly organizes your session into 4 panels
each set up for you to do certain tasks: Edit and write code (source
panel), run and execute code (console panel), list objects that you have
in your environment and have a history of your past commands
(environment/history panel), and a panel to show your folders'
structure, see plots, install/load packages, and read the documentation
of functions (files/plots/packages/help panel).

\includegraphics[width=900px]{assets/img/RStudio}

\begin{itemize}
\tightlist
\item
  the main script section, for writing scripts {[}top left section{]}
\item
  the console tab, for typing R commands directly {[}bottom left section
  as tab{]}
\item
  the terminal tab, for direct access to your system shell {[}bottom
  left section{]}
\item
  the plot tab, where you see the last plot generated {[}bottom right
  section as tab{]}
\item
  the help tab, with useful documentation of R functions {[}bottom right
  section as tab{]}
\item
  the history tab, with the list of the R commands used {[}top right
  section as tab{]}
\item
  the environment tab, with the created variables and functions loaded
  {[}top right section as tab{]}
\item
  and the packages tab, with the available/loaded R packages {[}bottom
  right section as tab{]}
\end{itemize}

\subsection{First steps with R}\label{first-steps-with-r}

This section is largely borrowed from the book Introduction to Data
Science by Rafael Irizarry (\url{https://rafalab.github.io/dsbook}).

\subsubsection{Objects}\label{objects}

Suppose a high school student asks us for help to solve several
quadratic equations of the form \(ax^2+bx+c = 0\). The quadratic formula
gives us the solutions:

\[
\frac{-b - \sqrt{b^2 - 4ac}}{2a}\,\, \mbox{ and } \frac{-b + \sqrt{b^2 - 4ac}}{2a}
\] which of course change depending on the values of \(a\), \(b\), and
\(c\). One advantage of programming languages is that we can define
variables and write expressions with these variables, similar to how we
do so in math, but obtain a numeric solution. We will write out general
code for the quadratic equation below, but if we are asked to solve
\(x^2 + x -1 = 0\), then we define:

\begin{Shaded}
\begin{Highlighting}[]
\NormalTok{a }\OtherTok{\textless{}{-}} \DecValTok{1}
\NormalTok{b }\OtherTok{\textless{}{-}} \DecValTok{1}
\NormalTok{c }\OtherTok{\textless{}{-}} \SpecialCharTok{{-}}\DecValTok{1}
\end{Highlighting}
\end{Shaded}

which stores the values for later use. We use \texttt{\textless{}-} to
assign values to the variables.

We can also assign values using \texttt{=} instead of
\texttt{\textless{}-}, but we recommend against using \texttt{=} to
avoid confusion.

Copy and paste the code above into your console to define the three
variables. Note that R does not print anything when we make this
assignment. This means the objects were defined successfully. Had you
made a mistake, you would have received an error message.

To see the value stored in a variable, we simply ask R to evaluate
\texttt{a} and it shows the stored value:

\begin{Shaded}
\begin{Highlighting}[]
\NormalTok{a}
\end{Highlighting}
\end{Shaded}

\begin{verbatim}
## [1] 1
\end{verbatim}

A more explicit way to ask R to show us the value stored in \texttt{a}
is using \texttt{print} like this:

\begin{Shaded}
\begin{Highlighting}[]
\FunctionTok{print}\NormalTok{(a)}
\end{Highlighting}
\end{Shaded}

\begin{verbatim}
## [1] 1
\end{verbatim}

We use the term \emph{object} to describe stuff that is stored in R.

\subsubsection{The workspace}\label{the-workspace}

As we define objects in the console, we are actually changing the
\emph{workspace}. You can see all the variables saved in your workspace
by typing:

\begin{Shaded}
\begin{Highlighting}[]
\FunctionTok{ls}\NormalTok{()}
\end{Highlighting}
\end{Shaded}

\begin{verbatim}
## [1] "a"    "b"    "c"    "path"
\end{verbatim}

In RStudio, the variables of the environment are displayed in the
\emph{Environment} tab.

We should see \texttt{a}, \texttt{b}, and \texttt{c}. If you try to
recover the value of a variable that is not in your workspace, you
receive an error. For example, if you type \texttt{x} you will receive
the following message:
\texttt{Error:\ object\ \textquotesingle{}x\textquotesingle{}\ not\ found}.

Now since these values are saved in variables, to obtain a solution to
our equation, we use the quadratic formula:

\begin{Shaded}
\begin{Highlighting}[]
\NormalTok{(}\SpecialCharTok{{-}}\NormalTok{b }\SpecialCharTok{+} \FunctionTok{sqrt}\NormalTok{(b}\SpecialCharTok{\^{}}\DecValTok{2} \SpecialCharTok{{-}} \DecValTok{4}\SpecialCharTok{*}\NormalTok{a}\SpecialCharTok{*}\NormalTok{c) ) }\SpecialCharTok{/}\NormalTok{ ( }\DecValTok{2}\SpecialCharTok{*}\NormalTok{a )}
\end{Highlighting}
\end{Shaded}

\begin{verbatim}
## [1] 0.618034
\end{verbatim}

\begin{Shaded}
\begin{Highlighting}[]
\NormalTok{(}\SpecialCharTok{{-}}\NormalTok{b }\SpecialCharTok{{-}} \FunctionTok{sqrt}\NormalTok{(b}\SpecialCharTok{\^{}}\DecValTok{2} \SpecialCharTok{{-}} \DecValTok{4}\SpecialCharTok{*}\NormalTok{a}\SpecialCharTok{*}\NormalTok{c) ) }\SpecialCharTok{/}\NormalTok{ ( }\DecValTok{2}\SpecialCharTok{*}\NormalTok{a )}
\end{Highlighting}
\end{Shaded}

\begin{verbatim}
## [1] -1.618034
\end{verbatim}

\subsubsection{Functions}\label{functions}

Once you define variables, the data analysis process can usually be
described as a series of \emph{functions} applied to the data. R
includes several predefined functions and most of the analysis pipelines
we construct make extensive use of these.

We already used the \texttt{print}, and \texttt{ls} functions. We also
used the function \texttt{sqrt} to solve the quadratic equation above.
There are many more prebuilt functions and even more can be added
through packages. These functions do not appear in the workspace because
you did not define them, but they are available for immediate use.

In general, we need to use parentheses to evaluate a function. If you
type \texttt{ls}, the function is not evaluated and instead R shows you
the code that defines the function. If you type \texttt{ls()} the
function is evaluated and, as seen above, we see objects in the
workspace.

Unlike \texttt{ls}, most functions require one or more \emph{arguments}.
Below is an example of how we assign an object to the argument of the
function \texttt{log}. Remember that we earlier defined \texttt{a} to be
1:

\begin{Shaded}
\begin{Highlighting}[]
\FunctionTok{log}\NormalTok{(}\DecValTok{8}\NormalTok{)}
\end{Highlighting}
\end{Shaded}

\begin{verbatim}
## [1] 2.079442
\end{verbatim}

\begin{Shaded}
\begin{Highlighting}[]
\FunctionTok{log}\NormalTok{(a) }
\end{Highlighting}
\end{Shaded}

\begin{verbatim}
## [1] 0
\end{verbatim}

You can find out what the function expects and what it does by reviewing
the very useful manuals included in R. You can get help by using the
\texttt{help} function like this:

\begin{Shaded}
\begin{Highlighting}[]
\FunctionTok{help}\NormalTok{(}\StringTok{"log"}\NormalTok{)}
\end{Highlighting}
\end{Shaded}

For most functions, we can also use this shorthand:

\begin{Shaded}
\begin{Highlighting}[]
\NormalTok{?log}
\end{Highlighting}
\end{Shaded}

The help page will show you what arguments the function is expecting.
For example, \texttt{log} needs \texttt{x} and \texttt{base} to run.
Some arguments are required and others are optional. You can determine
which arguments are optional by noting in the help document that a
default value is assigned with \texttt{=}. Defining these is optional.
For example, the base of the function \texttt{log} defaults to
\texttt{base\ =\ exp(1)} making \texttt{log} the natural log by default.

If you want a quick look at the arguments without opening the help
system, you can type:

\begin{Shaded}
\begin{Highlighting}[]
\FunctionTok{args}\NormalTok{(log)}
\end{Highlighting}
\end{Shaded}

\begin{verbatim}
## function (x, base = exp(1)) 
## NULL
\end{verbatim}

You can change the default values:

\begin{Shaded}
\begin{Highlighting}[]
\FunctionTok{log}\NormalTok{(}\DecValTok{8}\NormalTok{, }\AttributeTok{base =} \DecValTok{2}\NormalTok{)}
\end{Highlighting}
\end{Shaded}

\begin{verbatim}
## [1] 3
\end{verbatim}

Note that we have not been specifying the argument \texttt{x} as such:

\begin{Shaded}
\begin{Highlighting}[]
\FunctionTok{log}\NormalTok{(}\AttributeTok{x =} \DecValTok{8}\NormalTok{, }\AttributeTok{base =} \DecValTok{2}\NormalTok{)}
\end{Highlighting}
\end{Shaded}

\begin{verbatim}
## [1] 3
\end{verbatim}

The above code works, but we can save ourselves some typing: if no
argument name is used, R assumes you are entering arguments in the order
shown in the help file or by \texttt{args}. So by not using the names,
it assumes the arguments are \texttt{x} followed by \texttt{base}:

\begin{Shaded}
\begin{Highlighting}[]
\FunctionTok{log}\NormalTok{(}\DecValTok{8}\NormalTok{,}\DecValTok{2}\NormalTok{)}
\end{Highlighting}
\end{Shaded}

\begin{verbatim}
## [1] 3
\end{verbatim}

If we use the arguments' names, then we can include them in whatever
order we want:

\begin{Shaded}
\begin{Highlighting}[]
\FunctionTok{log}\NormalTok{(}\AttributeTok{base =} \DecValTok{2}\NormalTok{, }\AttributeTok{x =} \DecValTok{8}\NormalTok{)}
\end{Highlighting}
\end{Shaded}

\begin{verbatim}
## [1] 3
\end{verbatim}

To specify arguments, we must use \texttt{=}, and cannot use
\texttt{\textless{}-}.

There are some exceptions to the rule that functions need the
parentheses to be evaluated. Among these, the most commonly used are the
arithmetic and relational operators. For example:

\begin{Shaded}
\begin{Highlighting}[]
\DecValTok{2} \SpecialCharTok{\^{}} \DecValTok{3}
\end{Highlighting}
\end{Shaded}

\begin{verbatim}
## [1] 8
\end{verbatim}

You can see the arithmetic operators by typing:

\begin{Shaded}
\begin{Highlighting}[]
\FunctionTok{help}\NormalTok{(}\StringTok{"+"}\NormalTok{) }
\end{Highlighting}
\end{Shaded}

or

\begin{Shaded}
\begin{Highlighting}[]
\NormalTok{?}\StringTok{"+"}
\end{Highlighting}
\end{Shaded}

and the relational operators by typing:

\begin{Shaded}
\begin{Highlighting}[]
\FunctionTok{help}\NormalTok{(}\StringTok{"\textgreater{}"}\NormalTok{) }
\end{Highlighting}
\end{Shaded}

or

\begin{Shaded}
\begin{Highlighting}[]
\NormalTok{?}\StringTok{"\textgreater{}"}
\end{Highlighting}
\end{Shaded}

\subsubsection{Other prebuilt objects}\label{other-prebuilt-objects}

There are several datasets that are included for users to practice and
test out functions. You can see all the available datasets by typing:

\begin{Shaded}
\begin{Highlighting}[]
\FunctionTok{data}\NormalTok{()}
\end{Highlighting}
\end{Shaded}

This shows you the object name for these datasets. These datasets are
objects that can be used by simply typing the name. For example, if you
type:

\begin{Shaded}
\begin{Highlighting}[]
\NormalTok{co2}
\end{Highlighting}
\end{Shaded}

R will show you Mauna Loa atmospheric CO2 concentration data.

Other prebuilt objects are mathematical quantities, such as the constant
\(\pi\) and \(\infty\):

\begin{Shaded}
\begin{Highlighting}[]
\NormalTok{pi}
\end{Highlighting}
\end{Shaded}

\begin{verbatim}
## [1] 3.141593
\end{verbatim}

\begin{Shaded}
\begin{Highlighting}[]
\ConstantTok{Inf}\SpecialCharTok{+}\DecValTok{1}
\end{Highlighting}
\end{Shaded}

\begin{verbatim}
## [1] Inf
\end{verbatim}

\subsubsection{Variable names}\label{variable-names}

We have used the letters \texttt{a}, \texttt{b}, and \texttt{c} as
variable names, but variable names can be almost anything. Some basic
rules in R are that variable names have to start with a letter, can't
contain spaces, and should not be variables that are predefined in R.
For example, don't name one of your variables \texttt{install.packages}
by typing something like \texttt{install.packages\ \textless{}-\ 2}.

A nice convention to follow is to use meaningful words that describe
what is stored, use only lower case, and use underscores as a substitute
for spaces. For the quadratic equations, we could use something like
this:

\begin{Shaded}
\begin{Highlighting}[]
\NormalTok{solution\_1 }\OtherTok{\textless{}{-}}\NormalTok{ (}\SpecialCharTok{{-}}\NormalTok{b }\SpecialCharTok{+} \FunctionTok{sqrt}\NormalTok{(b}\SpecialCharTok{\^{}}\DecValTok{2} \SpecialCharTok{{-}} \DecValTok{4}\SpecialCharTok{*}\NormalTok{a}\SpecialCharTok{*}\NormalTok{c)) }\SpecialCharTok{/}\NormalTok{ (}\DecValTok{2}\SpecialCharTok{*}\NormalTok{a)}
\NormalTok{solution\_2 }\OtherTok{\textless{}{-}}\NormalTok{ (}\SpecialCharTok{{-}}\NormalTok{b }\SpecialCharTok{{-}} \FunctionTok{sqrt}\NormalTok{(b}\SpecialCharTok{\^{}}\DecValTok{2} \SpecialCharTok{{-}} \DecValTok{4}\SpecialCharTok{*}\NormalTok{a}\SpecialCharTok{*}\NormalTok{c)) }\SpecialCharTok{/}\NormalTok{ (}\DecValTok{2}\SpecialCharTok{*}\NormalTok{a)}
\end{Highlighting}
\end{Shaded}

For more advice, we highly recommend studying Hadley Wickham's style
guide\footnote{\url{http://adv-r.had.co.nz/Style.html}}.

\subsubsection{Reusing scripts}\label{reusing-scripts}

To solve another equation such as \(3x^2 + 2x -1\), we can copy and
paste the code above and then redefine the variables and recompute the
solution:

\begin{Shaded}
\begin{Highlighting}[]
\NormalTok{a }\OtherTok{\textless{}{-}} \DecValTok{3}
\NormalTok{b }\OtherTok{\textless{}{-}} \DecValTok{2}
\NormalTok{c }\OtherTok{\textless{}{-}} \SpecialCharTok{{-}}\DecValTok{1}
\NormalTok{(}\SpecialCharTok{{-}}\NormalTok{b }\SpecialCharTok{+} \FunctionTok{sqrt}\NormalTok{(b}\SpecialCharTok{\^{}}\DecValTok{2} \SpecialCharTok{{-}} \DecValTok{4}\SpecialCharTok{*}\NormalTok{a}\SpecialCharTok{*}\NormalTok{c)) }\SpecialCharTok{/}\NormalTok{ (}\DecValTok{2}\SpecialCharTok{*}\NormalTok{a)}
\NormalTok{(}\SpecialCharTok{{-}}\NormalTok{b }\SpecialCharTok{{-}} \FunctionTok{sqrt}\NormalTok{(b}\SpecialCharTok{\^{}}\DecValTok{2} \SpecialCharTok{{-}} \DecValTok{4}\SpecialCharTok{*}\NormalTok{a}\SpecialCharTok{*}\NormalTok{c)) }\SpecialCharTok{/}\NormalTok{ (}\DecValTok{2}\SpecialCharTok{*}\NormalTok{a)}
\end{Highlighting}
\end{Shaded}

By creating and saving a script with the code above, we would not need
to retype everything each time and, instead, simply change the variable
names. Try writing the script above into an editor and notice how easy
it is to change the variables and receive an answer.

\subsubsection{Commenting your code}\label{commenting-your-code}

If a line of R code starts with the symbol \texttt{\#}, it is not
evaluated. We can use this to write reminders of why we wrote particular
code. For example, in the script above we could add:

\begin{Shaded}
\begin{Highlighting}[]
\DocumentationTok{\#\# Code to compute solution to quadratic equation of the form ax\^{}2 + bx + c}
\DocumentationTok{\#\# define the variables}
\NormalTok{a }\OtherTok{\textless{}{-}} \DecValTok{3} 
\NormalTok{b }\OtherTok{\textless{}{-}} \DecValTok{2}
\NormalTok{c }\OtherTok{\textless{}{-}} \SpecialCharTok{{-}}\DecValTok{1}
\DocumentationTok{\#\# now compute the solution}
\NormalTok{(}\SpecialCharTok{{-}}\NormalTok{b }\SpecialCharTok{+} \FunctionTok{sqrt}\NormalTok{(b}\SpecialCharTok{\^{}}\DecValTok{2} \SpecialCharTok{{-}} \DecValTok{4}\SpecialCharTok{*}\NormalTok{a}\SpecialCharTok{*}\NormalTok{c)) }\SpecialCharTok{/}\NormalTok{ (}\DecValTok{2}\SpecialCharTok{*}\NormalTok{a)}
\NormalTok{(}\SpecialCharTok{{-}}\NormalTok{b }\SpecialCharTok{{-}} \FunctionTok{sqrt}\NormalTok{(b}\SpecialCharTok{\^{}}\DecValTok{2} \SpecialCharTok{{-}} \DecValTok{4}\SpecialCharTok{*}\NormalTok{a}\SpecialCharTok{*}\NormalTok{c)) }\SpecialCharTok{/}\NormalTok{ (}\DecValTok{2}\SpecialCharTok{*}\NormalTok{a)}
\end{Highlighting}
\end{Shaded}

\subsection{Data types}\label{data-types}

Variables in R can be of different types. For example, we need to
distinguish numbers from character strings and tables from simple lists
of numbers. The function \texttt{class} helps us determine what type of
object we have:

\begin{Shaded}
\begin{Highlighting}[]
\NormalTok{a }\OtherTok{\textless{}{-}} \DecValTok{2}
\FunctionTok{class}\NormalTok{(a)}
\end{Highlighting}
\end{Shaded}

\begin{verbatim}
## [1] "numeric"
\end{verbatim}

To work efficiently in R, it is important to learn the different types
of variables and what we can do with these.

\subsubsection{Data frames}\label{data-frames}

Up to now, the variables we have defined are just one number. We might
be interested in storing more complex data. The most common way of
storing a dataset in R is in a \emph{data frame}. Conceptually, we can
think of a data frame as a table with rows representing observations and
the different variables reported for each observation defining the
columns. Data frames are particularly useful for datasets because we can
combine different data types into one object.

A large proportion of data analysis challenges start with data stored in
a data frame.You can access this dataset by loading the \textbf{dslabs}
library and loading the \texttt{murders} dataset using the \texttt{data}
function:

\begin{Shaded}
\begin{Highlighting}[]
\FunctionTok{library}\NormalTok{(dslabs)}
\FunctionTok{data}\NormalTok{(murders)}
\end{Highlighting}
\end{Shaded}

To see that this is in fact a data frame, we type:

\begin{Shaded}
\begin{Highlighting}[]
\FunctionTok{class}\NormalTok{(murders)}
\end{Highlighting}
\end{Shaded}

\begin{verbatim}
## [1] "data.frame"
\end{verbatim}

\subsubsection{Examining an object}\label{examining-an-object}

The function \texttt{str} is useful for finding out more about the
structure of an object:

\begin{Shaded}
\begin{Highlighting}[]
\FunctionTok{str}\NormalTok{(murders)}
\end{Highlighting}
\end{Shaded}

\begin{verbatim}
## 'data.frame':    51 obs. of  5 variables:
## $ state : chr "Alabama" "Alaska" "Arizona" "Arkansas" ...
## $ abb : chr "AL" "AK" "AZ" "AR" ...
## $ region : Factor w/ 4 levels "Northeast","South",..: 2 4 4 2 4 4 1 2 2 2 ...
## $ population: num 4779736 710231 6392017 2915918 37253956 ...
## $ total : num 135 19 232 93 1257 ...
\end{verbatim}

This tells us much more about the object. We see that the table has 51
rows (50 states plus DC) and five variables (stored as columns). We can
show the first six lines using the function \texttt{head}:

\begin{Shaded}
\begin{Highlighting}[]
\FunctionTok{head}\NormalTok{(murders)  }\CommentTok{\# same as murders[1:6,]}
\end{Highlighting}
\end{Shaded}

\begin{verbatim}
##        state abb region population total
## 1    Alabama  AL  South    4779736   135
## 2     Alaska  AK   West     710231    19
## 3    Arizona  AZ   West    6392017   232
## 4   Arkansas  AR  South    2915918    93
## 5 California  CA   West   37253956  1257
## 6   Colorado  CO   West    5029196    65
\end{verbatim}

In this dataset, each state is considered an observation and five
variables are reported for each state.

Before we go any further in answering our original question about
different states, let's learn more about the components of this object.

\subsubsection{\texorpdfstring{The accessor:
\texttt{\$}}{The accessor: \$}}\label{the-accessor}

For our analysis, we will need to access the different variables
represented by columns included in this data frame. To do this, we use
the accessor operator \texttt{\$} in the following way:

\begin{Shaded}
\begin{Highlighting}[]
\NormalTok{murders}\SpecialCharTok{$}\NormalTok{population}
\end{Highlighting}
\end{Shaded}

\begin{verbatim}
##  [1]  4779736   710231  6392017  2915918 37253956  5029196  3574097   897934
##  [9]   601723 19687653  9920000  1360301  1567582 12830632  6483802  3046355
## [17]  2853118  4339367  4533372  1328361  5773552  6547629  9883640  5303925
## [25]  2967297  5988927   989415  1826341  2700551  1316470  8791894  2059179
## [33] 19378102  9535483   672591 11536504  3751351  3831074 12702379  1052567
## [41]  4625364   814180  6346105 25145561  2763885   625741  8001024  6724540
## [49]  1852994  5686986   563626
\end{verbatim}

But how did we know to use \texttt{population}? Previously, by applying
the function \texttt{str} to the object \texttt{murders}, we revealed
the names for each of the five variables stored in this table. We can
quickly access the variable names using:

\begin{Shaded}
\begin{Highlighting}[]
\FunctionTok{names}\NormalTok{(murders)}
\end{Highlighting}
\end{Shaded}

\begin{verbatim}
## [1] "state"      "abb"        "region"     "population" "total"
\end{verbatim}

It is important to know that the order of the entries in
\texttt{murders\$population} preserves the order of the rows in our data
table.

\textbf{Tip}: R comes with a very nice auto-complete functionality that
saves us the trouble of typing out all the names. Try typing
\texttt{murders\$p} then hitting the \emph{tab} key on your keyboard.
This functionality and many other useful auto-complete features are
available when working in RStudio.

\subsubsection{Vectors: numerics, characters, and
logical}\label{vectors-numerics-characters-and-logical}

The object \texttt{murders\$population} is not one number but several.
We call these types of objects \emph{vectors}. A single number is
technically a vector of length 1, but in general we use the term vectors
to refer to objects with several entries. The function \texttt{length}
tells you how many entries are in the vector:

\begin{Shaded}
\begin{Highlighting}[]
\NormalTok{pop }\OtherTok{\textless{}{-}}\NormalTok{ murders}\SpecialCharTok{$}\NormalTok{population}
\FunctionTok{length}\NormalTok{(pop)}
\end{Highlighting}
\end{Shaded}

\begin{verbatim}
## [1] 51
\end{verbatim}

This particular vector is \emph{numeric} since population sizes are
numbers:

\begin{Shaded}
\begin{Highlighting}[]
\FunctionTok{class}\NormalTok{(pop)}
\end{Highlighting}
\end{Shaded}

\begin{verbatim}
## [1] "numeric"
\end{verbatim}

In a numeric vector, every entry must be a number.

To store character strings, vectors can also be of class
\emph{character}. For example, the state names are characters:

\begin{Shaded}
\begin{Highlighting}[]
\FunctionTok{class}\NormalTok{(murders}\SpecialCharTok{$}\NormalTok{state)}
\end{Highlighting}
\end{Shaded}

\begin{verbatim}
## [1] "character"
\end{verbatim}

As with numeric vectors, all entries in a character vector need to be a
character.

Another important type of vectors are \emph{logical vectors}. These must
be either \texttt{TRUE} or \texttt{FALSE}.

\begin{Shaded}
\begin{Highlighting}[]
\NormalTok{z }\OtherTok{\textless{}{-}} \DecValTok{3} \SpecialCharTok{==} \DecValTok{2}
\NormalTok{z}
\end{Highlighting}
\end{Shaded}

\begin{verbatim}
## [1] FALSE
\end{verbatim}

\begin{Shaded}
\begin{Highlighting}[]
\FunctionTok{class}\NormalTok{(z)}
\end{Highlighting}
\end{Shaded}

\begin{verbatim}
## [1] "logical"
\end{verbatim}

Here the \texttt{==} is a relational operator asking if 3 is equal to 2.
In R, if you just use one \texttt{=}, you actually assign a variable,
but if you use two \texttt{==} you test for equality.

You can see the other \emph{relational operators} by typing:

\begin{Shaded}
\begin{Highlighting}[]
\NormalTok{?Comparison}
\end{Highlighting}
\end{Shaded}

In future sections, you will see how useful relational operators can be.

\textbf{Advanced}: Mathematically, the values in \texttt{pop} are
integers and there is an integer class in R. However, by default,
numbers are assigned class numeric even when they are round integers.
For example, \texttt{class(1)} returns numeric. You can turn them into
class integer with the \texttt{as.integer()} function or by adding an
\texttt{L} like this: \texttt{1L}. Note the class by typing:
\texttt{class(1L)}

\subsubsection{Factors}\label{factors}

In the \texttt{murders} dataset, we might expect the region to also be a
character vector. However, it is not:

\begin{Shaded}
\begin{Highlighting}[]
\FunctionTok{class}\NormalTok{(murders}\SpecialCharTok{$}\NormalTok{region)}
\end{Highlighting}
\end{Shaded}

\begin{verbatim}
## [1] "factor"
\end{verbatim}

It is a \emph{factor}. Factors are useful for storing categorical data.
We can see that there are only 4 regions by using the \texttt{levels}
function:

\begin{Shaded}
\begin{Highlighting}[]
\FunctionTok{levels}\NormalTok{(murders}\SpecialCharTok{$}\NormalTok{region)}
\end{Highlighting}
\end{Shaded}

\begin{verbatim}
## [1] "Northeast"     "South"         "North Central" "West"
\end{verbatim}

In the background, R stores these \emph{levels} as integers and keeps a
map to keep track of the labels. This is more memory efficient than
storing all the characters.

Note that the levels have an order that is different from the order of
appearance in the factor object. The default in R is for the levels to
follow alphabetical order. However, often we want the levels to follow a
different order. You can specify an order through the \texttt{levels}
argument when creating the factor with the \texttt{factor} function. For
example, in the murders dataset regions are ordered from east to west.
The function \texttt{reorder} lets us change the order of the levels of
a factor variable based on a summary computed on a numeric vector. We
will demonstrate this with a simple example, and will see more advanced
ones in the Data Visualization part of the book.

Suppose we want the levels of the region by the total number of murders
rather than alphabetical order. If there are values associated with each
level, we can use the \texttt{reorder} and specify a data summary to
determine the order. The following code takes the sum of the total
murders in each region, and reorders the factor following these sums.

\begin{Shaded}
\begin{Highlighting}[]
\NormalTok{region }\OtherTok{\textless{}{-}}\NormalTok{ murders}\SpecialCharTok{$}\NormalTok{region}
\NormalTok{value }\OtherTok{\textless{}{-}}\NormalTok{ murders}\SpecialCharTok{$}\NormalTok{total}
\NormalTok{region }\OtherTok{\textless{}{-}} \FunctionTok{reorder}\NormalTok{(region, value, }\AttributeTok{FUN =}\NormalTok{ sum)}
\FunctionTok{levels}\NormalTok{(region)}
\end{Highlighting}
\end{Shaded}

\begin{verbatim}
## [1] "Northeast"     "North Central" "West"          "South"
\end{verbatim}

The new order is in agreement with the fact that the Northeast has the
least murders and the South has the most.

\textbf{Warning}: Factors can be a source of confusion since sometimes
they behave like characters and sometimes they do not. As a result,
confusing factors and characters are a common source of bugs.

\subsubsection{Lists}\label{lists}

Data frames are a special case of \emph{lists}. Lists are useful because
they can store any combination of different types. You can create a list
using the \texttt{list} function like this:

\begin{Shaded}
\begin{Highlighting}[]
\NormalTok{record }\OtherTok{\textless{}{-}} \FunctionTok{list}\NormalTok{(}\AttributeTok{name =} \StringTok{"John Doe"}\NormalTok{,}
             \AttributeTok{student\_id =} \DecValTok{1234}\NormalTok{,}
             \AttributeTok{grades =} \FunctionTok{c}\NormalTok{(}\DecValTok{95}\NormalTok{, }\DecValTok{82}\NormalTok{, }\DecValTok{91}\NormalTok{, }\DecValTok{97}\NormalTok{, }\DecValTok{93}\NormalTok{),}
             \AttributeTok{final\_grade =} \StringTok{"A"}\NormalTok{)}
\end{Highlighting}
\end{Shaded}

The function \texttt{c} is described in Section @ref(vectors).

This list includes a character, a number, a vector with five numbers,
and another character.

\begin{Shaded}
\begin{Highlighting}[]
\NormalTok{record}
\end{Highlighting}
\end{Shaded}

\begin{verbatim}
## $name
## [1] "John Doe"
## 
## $student_id
## [1] 1234
## 
## $grades
## [1] 95 82 91 97 93
## 
## $final_grade
## [1] "A"
\end{verbatim}

\begin{Shaded}
\begin{Highlighting}[]
\FunctionTok{class}\NormalTok{(record)}
\end{Highlighting}
\end{Shaded}

\begin{verbatim}
## [1] "list"
\end{verbatim}

As with data frames, you can extract the components of a list with the
accessor \texttt{\$}.

\begin{Shaded}
\begin{Highlighting}[]
\NormalTok{record}\SpecialCharTok{$}\NormalTok{student\_id}
\end{Highlighting}
\end{Shaded}

\begin{verbatim}
## [1] 1234
\end{verbatim}

We can also use double square brackets (\texttt{{[}{[}}) like this:

\begin{Shaded}
\begin{Highlighting}[]
\NormalTok{record[[}\StringTok{"student\_id"}\NormalTok{]]}
\end{Highlighting}
\end{Shaded}

\begin{verbatim}
## [1] 1234
\end{verbatim}

You should get used to the fact that in R, there are often several ways
to do the same thing, such as accessing entries.

You might also encounter lists without variable names.

\begin{Shaded}
\begin{Highlighting}[]
\NormalTok{record2 }\OtherTok{\textless{}{-}} \FunctionTok{list}\NormalTok{(}\StringTok{"John Doe"}\NormalTok{, }\DecValTok{1234}\NormalTok{)}
\NormalTok{record2}
\end{Highlighting}
\end{Shaded}

\begin{verbatim}
## [[1]]
## [1] "John Doe"
## 
## [[2]]
## [1] 1234
\end{verbatim}

If a list does not have names, you cannot extract the elements with
\texttt{\$}, but you can still use the brackets method and instead of
providing the variable name, you provide the list index, like this:

\begin{Shaded}
\begin{Highlighting}[]
\NormalTok{record2[[}\DecValTok{1}\NormalTok{]]}
\end{Highlighting}
\end{Shaded}

\begin{verbatim}
## [1] "John Doe"
\end{verbatim}

\subsubsection{Matrices}\label{matrices}

Matrices are another very common object type in R. Matrices are similar
to data frames in that they are two-dimensional: they have rows and
columns. However, like numeric, character and logical vectors, entries
in matrices have to be all the same type. For this reason data frames
are much more useful for storing data, since we can have characters,
factors, and numbers in them.

Yet matrices have a major advantage over data frames: we can perform
matrix algebra operations, a powerful type of mathematical technique. We
do not describe these operations in this book, but much of what happens
in the background when you perform a data analysis involves matrices.

We can define a matrix using the \texttt{matrix} function. We need to
specify the number of rows and columns.

\begin{Shaded}
\begin{Highlighting}[]
\NormalTok{mat }\OtherTok{\textless{}{-}} \FunctionTok{matrix}\NormalTok{(}\DecValTok{1}\SpecialCharTok{:}\DecValTok{12}\NormalTok{, }\AttributeTok{nrow =} \DecValTok{4}\NormalTok{, }\AttributeTok{ncol =} \DecValTok{3}\NormalTok{)}
\NormalTok{mat}
\end{Highlighting}
\end{Shaded}

\begin{verbatim}
##      [,1] [,2] [,3]
## [1,]    1    5    9
## [2,]    2    6   10
## [3,]    3    7   11
## [4,]    4    8   12
\end{verbatim}

You can access specific entries in a matrix using square brackets
(\texttt{{[}}). If you want the second row, third column, you use:

\begin{Shaded}
\begin{Highlighting}[]
\NormalTok{mat[}\DecValTok{2}\NormalTok{, }\DecValTok{3}\NormalTok{]}
\end{Highlighting}
\end{Shaded}

\begin{verbatim}
## [1] 10
\end{verbatim}

If you want the entire second row, you leave the column spot empty:

\begin{Shaded}
\begin{Highlighting}[]
\NormalTok{mat[}\DecValTok{2}\NormalTok{, ]}
\end{Highlighting}
\end{Shaded}

\begin{verbatim}
## [1]  2  6 10
\end{verbatim}

Notice that this returns a vector, not a matrix.

Similarly, if you want the entire third column, you leave the row spot
empty:

\begin{Shaded}
\begin{Highlighting}[]
\NormalTok{mat[, }\DecValTok{3}\NormalTok{]}
\end{Highlighting}
\end{Shaded}

\begin{verbatim}
## [1]  9 10 11 12
\end{verbatim}

This is also a vector, not a matrix.

You can access more than one column or more than one row if you like.
This will give you a new matrix.

\begin{Shaded}
\begin{Highlighting}[]
\NormalTok{mat[, }\DecValTok{2}\SpecialCharTok{:}\DecValTok{3}\NormalTok{]}
\end{Highlighting}
\end{Shaded}

\begin{verbatim}
##      [,1] [,2]
## [1,]    5    9
## [2,]    6   10
## [3,]    7   11
## [4,]    8   12
\end{verbatim}

You can subset both rows and columns:

\begin{Shaded}
\begin{Highlighting}[]
\NormalTok{mat[}\DecValTok{1}\SpecialCharTok{:}\DecValTok{2}\NormalTok{, }\DecValTok{2}\SpecialCharTok{:}\DecValTok{3}\NormalTok{]}
\end{Highlighting}
\end{Shaded}

\begin{verbatim}
##      [,1] [,2]
## [1,]    5    9
## [2,]    6   10
\end{verbatim}

We can convert matrices into data frames using the function
\texttt{as.data.frame}:

\begin{Shaded}
\begin{Highlighting}[]
\FunctionTok{as.data.frame}\NormalTok{(mat)}
\end{Highlighting}
\end{Shaded}

\begin{verbatim}
##   V1 V2 V3
## 1  1  5  9
## 2  2  6 10
## 3  3  7 11
## 4  4  8 12
\end{verbatim}

You can also use single square brackets (\texttt{{[}}) to access rows
and columns of a data frame:

\begin{Shaded}
\begin{Highlighting}[]
\FunctionTok{data}\NormalTok{(}\StringTok{"murders"}\NormalTok{)}
\NormalTok{murders[}\DecValTok{25}\NormalTok{, }\DecValTok{1}\NormalTok{]}
\end{Highlighting}
\end{Shaded}

\begin{verbatim}
## [1] "Mississippi"
\end{verbatim}

\begin{Shaded}
\begin{Highlighting}[]
\NormalTok{murders[}\DecValTok{2}\SpecialCharTok{:}\DecValTok{3}\NormalTok{, ]}
\end{Highlighting}
\end{Shaded}

\begin{verbatim}
##     state abb region population total
## 2  Alaska  AK   West     710231    19
## 3 Arizona  AZ   West    6392017   232
\end{verbatim}

\subsection{Vectors}\label{vectors}

In R, the most basic objects available to store data are \emph{vectors}.
As we have seen, complex datasets can usually be broken down into
components that are vectors. For example, in a data frame, each column
is a vector. Here we learn more about this important class.

\subsubsection{Creating vectors}\label{creating-vectors}

We can create vectors using the function \texttt{c}, which stands for
\emph{concatenate}. We use \texttt{c} to concatenate entries in the
following way:

\begin{Shaded}
\begin{Highlighting}[]
\NormalTok{codes }\OtherTok{\textless{}{-}} \FunctionTok{c}\NormalTok{(}\DecValTok{380}\NormalTok{, }\DecValTok{124}\NormalTok{, }\DecValTok{818}\NormalTok{)}
\NormalTok{codes}
\end{Highlighting}
\end{Shaded}

\begin{verbatim}
## [1] 380 124 818
\end{verbatim}

We can also create character vectors. We use the quotes to denote that
the entries are characters rather than variable names.

\begin{Shaded}
\begin{Highlighting}[]
\NormalTok{country }\OtherTok{\textless{}{-}} \FunctionTok{c}\NormalTok{(}\StringTok{"italy"}\NormalTok{, }\StringTok{"canada"}\NormalTok{, }\StringTok{"egypt"}\NormalTok{)}
\end{Highlighting}
\end{Shaded}

In R you can also use single quotes:

\begin{Shaded}
\begin{Highlighting}[]
\NormalTok{country }\OtherTok{\textless{}{-}} \FunctionTok{c}\NormalTok{(}\StringTok{\textquotesingle{}italy\textquotesingle{}}\NormalTok{, }\StringTok{\textquotesingle{}canada\textquotesingle{}}\NormalTok{, }\StringTok{\textquotesingle{}egypt\textquotesingle{}}\NormalTok{)}
\end{Highlighting}
\end{Shaded}

But be careful not to confuse the single quote ' with the \emph{back
quote} `. By now you should know that if you type:

\begin{Shaded}
\begin{Highlighting}[]
\NormalTok{country }\OtherTok{\textless{}{-}} \FunctionTok{c}\NormalTok{(italy, canada, egypt)}
\end{Highlighting}
\end{Shaded}

you receive an error because the variables \texttt{italy},
\texttt{canada}, and \texttt{egypt} are not defined. If we do not use
the quotes, R looks for variables with those names and returns an error.

\subsubsection{Names}\label{names}

Sometimes it is useful to name the entries of a vector. For example,
when defining a vector of country codes, we can use the names to connect
the two:

\begin{Shaded}
\begin{Highlighting}[]
\NormalTok{codes }\OtherTok{\textless{}{-}} \FunctionTok{c}\NormalTok{(}\AttributeTok{italy =} \DecValTok{380}\NormalTok{, }\AttributeTok{canada =} \DecValTok{124}\NormalTok{, }\AttributeTok{egypt =} \DecValTok{818}\NormalTok{)}
\NormalTok{codes}
\end{Highlighting}
\end{Shaded}

\begin{verbatim}
##  italy canada  egypt 
##    380    124    818
\end{verbatim}

The object \texttt{codes} continues to be a numeric vector:

\begin{Shaded}
\begin{Highlighting}[]
\FunctionTok{class}\NormalTok{(codes)}
\end{Highlighting}
\end{Shaded}

\begin{verbatim}
## [1] "numeric"
\end{verbatim}

but with names:

\begin{Shaded}
\begin{Highlighting}[]
\FunctionTok{names}\NormalTok{(codes)}
\end{Highlighting}
\end{Shaded}

\begin{verbatim}
## [1] "italy"  "canada" "egypt"
\end{verbatim}

If the use of strings without quotes looks confusing, know that you can
use the quotes as well:

\begin{Shaded}
\begin{Highlighting}[]
\NormalTok{codes }\OtherTok{\textless{}{-}} \FunctionTok{c}\NormalTok{(}\StringTok{"italy"} \OtherTok{=} \DecValTok{380}\NormalTok{, }\StringTok{"canada"} \OtherTok{=} \DecValTok{124}\NormalTok{, }\StringTok{"egypt"} \OtherTok{=} \DecValTok{818}\NormalTok{)}
\NormalTok{codes}
\end{Highlighting}
\end{Shaded}

\begin{verbatim}
##  italy canada  egypt 
##    380    124    818
\end{verbatim}

There is no difference between this function call and the previous one.
This is one of the many ways in which R is quirky compared to other
languages.

We can also assign names using the \texttt{names} functions:

\begin{Shaded}
\begin{Highlighting}[]
\NormalTok{codes }\OtherTok{\textless{}{-}} \FunctionTok{c}\NormalTok{(}\DecValTok{380}\NormalTok{, }\DecValTok{124}\NormalTok{, }\DecValTok{818}\NormalTok{)}
\NormalTok{country }\OtherTok{\textless{}{-}} \FunctionTok{c}\NormalTok{(}\StringTok{"italy"}\NormalTok{,}\StringTok{"canada"}\NormalTok{,}\StringTok{"egypt"}\NormalTok{)}
\FunctionTok{names}\NormalTok{(codes) }\OtherTok{\textless{}{-}}\NormalTok{ country}
\NormalTok{codes}
\end{Highlighting}
\end{Shaded}

\begin{verbatim}
##  italy canada  egypt 
##    380    124    818
\end{verbatim}

\subsubsection{Sequences}\label{sequences}

Another useful function for creating vectors in sequences is
\texttt{seq}:

\begin{Shaded}
\begin{Highlighting}[]
\FunctionTok{seq}\NormalTok{(}\DecValTok{1}\NormalTok{, }\DecValTok{10}\NormalTok{)}
\end{Highlighting}
\end{Shaded}

\begin{verbatim}
##  [1]  1  2  3  4  5  6  7  8  9 10
\end{verbatim}

The first argument defines the start, and the second defines the end
which is included. The default is to go up in increments of 1, but a
third argument lets us tell it how much to jump by:

\begin{Shaded}
\begin{Highlighting}[]
\FunctionTok{seq}\NormalTok{(}\DecValTok{1}\NormalTok{, }\DecValTok{10}\NormalTok{, }\DecValTok{2}\NormalTok{)}
\end{Highlighting}
\end{Shaded}

\begin{verbatim}
## [1] 1 3 5 7 9
\end{verbatim}

If we want consecutive integers, we can use the following shorthand:

\begin{Shaded}
\begin{Highlighting}[]
\DecValTok{1}\SpecialCharTok{:}\DecValTok{10}
\end{Highlighting}
\end{Shaded}

\begin{verbatim}
##  [1]  1  2  3  4  5  6  7  8  9 10
\end{verbatim}

When we use these functions, R produces integers, not numerics, because
they are typically used to index something:

\begin{Shaded}
\begin{Highlighting}[]
\FunctionTok{class}\NormalTok{(}\DecValTok{1}\SpecialCharTok{:}\DecValTok{10}\NormalTok{)}
\end{Highlighting}
\end{Shaded}

\begin{verbatim}
## [1] "integer"
\end{verbatim}

However, if we create a sequence including non-integers, the class
changes:

\begin{Shaded}
\begin{Highlighting}[]
\FunctionTok{class}\NormalTok{(}\FunctionTok{seq}\NormalTok{(}\DecValTok{1}\NormalTok{, }\DecValTok{10}\NormalTok{, }\FloatTok{0.5}\NormalTok{))}
\end{Highlighting}
\end{Shaded}

\begin{verbatim}
## [1] "numeric"
\end{verbatim}

\subsubsection{Subsetting}\label{subsetting}

We use square brackets to access specific elements of a vector. For the
vector \texttt{codes} we defined above, we can access the second element
using:

\begin{Shaded}
\begin{Highlighting}[]
\NormalTok{codes[}\DecValTok{2}\NormalTok{]}
\end{Highlighting}
\end{Shaded}

\begin{verbatim}
## canada 
##    124
\end{verbatim}

You can get more than one entry by using a multi-entry vector as an
index:

\begin{Shaded}
\begin{Highlighting}[]
\NormalTok{codes[}\FunctionTok{c}\NormalTok{(}\DecValTok{1}\NormalTok{,}\DecValTok{3}\NormalTok{)]}
\end{Highlighting}
\end{Shaded}

\begin{verbatim}
## italy egypt 
##   380   818
\end{verbatim}

The sequences defined above are particularly useful if we want to
access, say, the first two elements:

\begin{Shaded}
\begin{Highlighting}[]
\NormalTok{codes[}\DecValTok{1}\SpecialCharTok{:}\DecValTok{2}\NormalTok{]}
\end{Highlighting}
\end{Shaded}

\begin{verbatim}
##  italy canada 
##    380    124
\end{verbatim}

If the elements have names, we can also access the entries using these
names. Below are two examples.

\begin{Shaded}
\begin{Highlighting}[]
\NormalTok{codes[}\StringTok{"canada"}\NormalTok{]}
\end{Highlighting}
\end{Shaded}

\begin{verbatim}
## canada 
##    124
\end{verbatim}

\begin{Shaded}
\begin{Highlighting}[]
\NormalTok{codes[}\FunctionTok{c}\NormalTok{(}\StringTok{"egypt"}\NormalTok{,}\StringTok{"italy"}\NormalTok{)]}
\end{Highlighting}
\end{Shaded}

\begin{verbatim}
## egypt italy 
##   818   380
\end{verbatim}

\subsection{Coercion}\label{coercion}

In general, \emph{coercion} is an attempt by R to be flexible with data
types. When an entry does not match the expected, some of the prebuilt R
functions try to guess what was meant before throwing an error. This can
also lead to confusion. Failing to understand \emph{coercion} can drive
programmers crazy when attempting to code in R since it behaves quite
differently from most other languages in this regard. Let's learn about
it with some examples.

We said that vectors must be all of the same type. So if we try to
combine, say, numbers and characters, you might expect an error:

\begin{Shaded}
\begin{Highlighting}[]
\NormalTok{x }\OtherTok{\textless{}{-}} \FunctionTok{c}\NormalTok{(}\DecValTok{1}\NormalTok{, }\StringTok{"canada"}\NormalTok{, }\DecValTok{3}\NormalTok{)}
\end{Highlighting}
\end{Shaded}

But we don't get one, not even a warning! What happened? Look at
\texttt{x} and its class:

\begin{Shaded}
\begin{Highlighting}[]
\NormalTok{x}
\end{Highlighting}
\end{Shaded}

\begin{verbatim}
## [1] "1"      "canada" "3"
\end{verbatim}

\begin{Shaded}
\begin{Highlighting}[]
\FunctionTok{class}\NormalTok{(x)}
\end{Highlighting}
\end{Shaded}

\begin{verbatim}
## [1] "character"
\end{verbatim}

R \emph{coerced} the data into characters. It guessed that because you
put a character string in the vector, you meant the 1 and 3 to actually
be character strings \texttt{"1"} and ``\texttt{3}''. The fact that not
even a warning is issued is an example of how coercion can cause many
unnoticed errors in R.

R also offers functions to change from one type to another. For example,
you can turn numbers into characters with:

\begin{Shaded}
\begin{Highlighting}[]
\NormalTok{x }\OtherTok{\textless{}{-}} \DecValTok{1}\SpecialCharTok{:}\DecValTok{5}
\NormalTok{y }\OtherTok{\textless{}{-}} \FunctionTok{as.character}\NormalTok{(x)}
\NormalTok{y}
\end{Highlighting}
\end{Shaded}

\begin{verbatim}
## [1] "1" "2" "3" "4" "5"
\end{verbatim}

You can turn it back with \texttt{as.numeric}:

\begin{Shaded}
\begin{Highlighting}[]
\FunctionTok{as.numeric}\NormalTok{(y)}
\end{Highlighting}
\end{Shaded}

\begin{verbatim}
## [1] 1 2 3 4 5
\end{verbatim}

This function is actually quite useful since datasets that include
numbers as character strings are common.

\subsubsection{Not availables (NA)}\label{not-availables-na}

When a function tries to coerce one type to another and encounters an
impossible case, it usually gives us a warning and turns the entry into
a special value called an \texttt{NA} for ``not available''. For
example:

\begin{Shaded}
\begin{Highlighting}[]
\NormalTok{x }\OtherTok{\textless{}{-}} \FunctionTok{c}\NormalTok{(}\StringTok{"1"}\NormalTok{, }\StringTok{"b"}\NormalTok{, }\StringTok{"3"}\NormalTok{)}
\FunctionTok{as.numeric}\NormalTok{(x)}
\end{Highlighting}
\end{Shaded}

\begin{verbatim}
## Warning: NAs introduced by coercion
\end{verbatim}

\begin{verbatim}
## [1]  1 NA  3
\end{verbatim}

R does not have any guesses for what number you want when you type
\texttt{b}, so it returns \texttt{NA}.

As a data scientist you will encounter the \texttt{NA}s often as they
are generally used for missing data, a common problem in real-world
datasets.

\subsection{Sorting}\label{sorting}

Now that we have mastered some basic R knowledge, let's try to gain some
insights into the safety of different states in the context of gun
murders.

\subsubsection{\texorpdfstring{\texttt{sort}}{sort}}\label{sort}

Say we want to rank the states from least to most gun murders. The
function \texttt{sort} sorts a vector in increasing order. We can
therefore see the largest number of gun murders by typing:

\begin{Shaded}
\begin{Highlighting}[]
\FunctionTok{library}\NormalTok{(dslabs)}
\FunctionTok{data}\NormalTok{(murders)}
\FunctionTok{sort}\NormalTok{(murders}\SpecialCharTok{$}\NormalTok{total)}
\end{Highlighting}
\end{Shaded}

\begin{verbatim}
##  [1]    2    4    5    5    7    8   11   12   12   16   19   21   22   27   32
## [16]   36   38   53   63   65   67   84   93   93   97   97   99  111  116  118
## [31]  120  135  142  207  219  232  246  250  286  293  310  321  351  364  376
## [46]  413  457  517  669  805 1257
\end{verbatim}

However, this does not give us information about which states have which
murder totals. For example, we don't know which state had 1257.

\subsubsection{\texorpdfstring{\texttt{order}}{order}}\label{order}

The function \texttt{order} is closer to what we want. It takes a vector
as input and returns the vector of indexes that sorts the input vector.
This may sound confusing so let's look at a simple example. We can
create a vector and sort it:

\begin{Shaded}
\begin{Highlighting}[]
\NormalTok{x }\OtherTok{\textless{}{-}} \FunctionTok{c}\NormalTok{(}\DecValTok{31}\NormalTok{, }\DecValTok{4}\NormalTok{, }\DecValTok{15}\NormalTok{, }\DecValTok{92}\NormalTok{, }\DecValTok{65}\NormalTok{)}
\FunctionTok{sort}\NormalTok{(x)}
\end{Highlighting}
\end{Shaded}

\begin{verbatim}
## [1]  4 15 31 65 92
\end{verbatim}

Rather than sorting the input vector, the function \texttt{order}
returns the indexes that sort the input vector:

\begin{Shaded}
\begin{Highlighting}[]
\NormalTok{index }\OtherTok{\textless{}{-}} \FunctionTok{order}\NormalTok{(x)}
\NormalTok{x[index]}
\end{Highlighting}
\end{Shaded}

\begin{verbatim}
## [1]  4 15 31 65 92
\end{verbatim}

This is the same output as that returned by \texttt{sort(x)}. If we look
at this index, we see why it works:

\begin{Shaded}
\begin{Highlighting}[]
\NormalTok{x}
\end{Highlighting}
\end{Shaded}

\begin{verbatim}
## [1] 31  4 15 92 65
\end{verbatim}

\begin{Shaded}
\begin{Highlighting}[]
\FunctionTok{order}\NormalTok{(x)}
\end{Highlighting}
\end{Shaded}

\begin{verbatim}
## [1] 2 3 1 5 4
\end{verbatim}

The second entry of \texttt{x} is the smallest, so \texttt{order(x)}
starts with \texttt{2}. The next smallest is the third entry, so the
second entry is \texttt{3} and so on.

How does this help us order the states by murders? First, remember that
the entries of vectors you access with \texttt{\$} follow the same order
as the rows in the table. For example, these two vectors containing
state names and abbreviations, respectively, are matched by their order:

\begin{Shaded}
\begin{Highlighting}[]
\NormalTok{murders}\SpecialCharTok{$}\NormalTok{state[}\DecValTok{1}\SpecialCharTok{:}\DecValTok{6}\NormalTok{]}
\end{Highlighting}
\end{Shaded}

\begin{verbatim}
## [1] "Alabama"    "Alaska"     "Arizona"    "Arkansas"   "California"
## [6] "Colorado"
\end{verbatim}

\begin{Shaded}
\begin{Highlighting}[]
\NormalTok{murders}\SpecialCharTok{$}\NormalTok{abb[}\DecValTok{1}\SpecialCharTok{:}\DecValTok{6}\NormalTok{]}
\end{Highlighting}
\end{Shaded}

\begin{verbatim}
## [1] "AL" "AK" "AZ" "AR" "CA" "CO"
\end{verbatim}

This means we can order the state names by their total murders. We first
obtain the index that orders the vectors according to murder totals and
then index the state names vector:

\begin{Shaded}
\begin{Highlighting}[]
\NormalTok{ind }\OtherTok{\textless{}{-}} \FunctionTok{order}\NormalTok{(murders}\SpecialCharTok{$}\NormalTok{total)}
\NormalTok{murders}\SpecialCharTok{$}\NormalTok{abb[ind] }
\end{Highlighting}
\end{Shaded}

\begin{verbatim}
##  [1] "VT" "ND" "NH" "WY" "HI" "SD" "ME" "ID" "MT" "RI" "AK" "IA" "UT" "WV" "NE"
## [16] "OR" "DE" "MN" "KS" "CO" "NM" "NV" "AR" "WA" "CT" "WI" "DC" "OK" "KY" "MA"
## [31] "MS" "AL" "IN" "SC" "TN" "AZ" "NJ" "VA" "NC" "MD" "OH" "MO" "LA" "IL" "GA"
## [46] "MI" "PA" "NY" "FL" "TX" "CA"
\end{verbatim}

According to the above, California (CA) had the most murders.

\subsubsection{\texorpdfstring{\texttt{max} and
\texttt{which.max}}{max and which.max}}\label{max-and-which.max}

If we are only interested in the entry with the largest value, we can
use \texttt{max} for the value:

\begin{Shaded}
\begin{Highlighting}[]
\FunctionTok{max}\NormalTok{(murders}\SpecialCharTok{$}\NormalTok{total)}
\end{Highlighting}
\end{Shaded}

\begin{verbatim}
## [1] 1257
\end{verbatim}

and \texttt{which.max} for the index of the largest value:

\begin{Shaded}
\begin{Highlighting}[]
\NormalTok{i\_max }\OtherTok{\textless{}{-}} \FunctionTok{which.max}\NormalTok{(murders}\SpecialCharTok{$}\NormalTok{total)}
\NormalTok{murders}\SpecialCharTok{$}\NormalTok{state[i\_max]}
\end{Highlighting}
\end{Shaded}

\begin{verbatim}
## [1] "California"
\end{verbatim}

For the minimum, we can use \texttt{min} and \texttt{which.min} in the
same way.

Does this mean California is the most dangerous state? In an upcoming
section, we argue that we should be considering rates instead of totals.
Before doing that, we introduce one last order-related function:
\texttt{rank}.

\subsubsection{\texorpdfstring{\texttt{rank}}{rank}}\label{rank}

Although not as frequently used as \texttt{order} and \texttt{sort}, the
function \texttt{rank} is also related to order and can be useful. For
any given vector it returns a vector with the rank of the first entry,
second entry, etc., of the input vector. Here is a simple example:

\begin{Shaded}
\begin{Highlighting}[]
\NormalTok{x }\OtherTok{\textless{}{-}} \FunctionTok{c}\NormalTok{(}\DecValTok{31}\NormalTok{, }\DecValTok{4}\NormalTok{, }\DecValTok{15}\NormalTok{, }\DecValTok{92}\NormalTok{, }\DecValTok{65}\NormalTok{)}
\FunctionTok{rank}\NormalTok{(x)}
\end{Highlighting}
\end{Shaded}

\begin{verbatim}
## [1] 3 1 2 5 4
\end{verbatim}

\subsubsection{Beware of recycling}\label{beware-of-recycling}

Another common source of unnoticed errors in R is the use of
\emph{recycling}. We saw that vectors are added element-wise. So if the
vectors don't match in length, it is natural to assume that we should
get an error. But we don't. Notice what happens:

\begin{Shaded}
\begin{Highlighting}[]
\NormalTok{x }\OtherTok{\textless{}{-}} \FunctionTok{c}\NormalTok{(}\DecValTok{1}\NormalTok{,}\DecValTok{2}\NormalTok{,}\DecValTok{3}\NormalTok{)}
\NormalTok{y }\OtherTok{\textless{}{-}} \FunctionTok{c}\NormalTok{(}\DecValTok{10}\NormalTok{, }\DecValTok{20}\NormalTok{, }\DecValTok{30}\NormalTok{, }\DecValTok{40}\NormalTok{, }\DecValTok{50}\NormalTok{, }\DecValTok{60}\NormalTok{, }\DecValTok{70}\NormalTok{)}
\NormalTok{x}\SpecialCharTok{+}\NormalTok{y}
\end{Highlighting}
\end{Shaded}

\begin{verbatim}
## Warning in x + y: longer object length is not a multiple of shorter object
## length
\end{verbatim}

\begin{verbatim}
## [1] 11 22 33 41 52 63 71
\end{verbatim}

We do get a warning, but no error. For the output, R has recycled the
numbers in \texttt{x}. Notice the last digit of numbers in the output.

\subsection{Vector arithmetics}\label{vector-arithmetics}

California had the most murders, but does this mean it is the most
dangerous state? What if it just has many more people than any other
state? We can quickly confirm that California indeed has the largest
population:

\begin{Shaded}
\begin{Highlighting}[]
\FunctionTok{library}\NormalTok{(dslabs)}
\FunctionTok{data}\NormalTok{(}\StringTok{"murders"}\NormalTok{)}
\NormalTok{murders}\SpecialCharTok{$}\NormalTok{state[}\FunctionTok{which.max}\NormalTok{(murders}\SpecialCharTok{$}\NormalTok{population)]}
\end{Highlighting}
\end{Shaded}

\begin{verbatim}
## [1] "California"
\end{verbatim}

with over 37 million inhabitants. It is therefore unfair to compare the
totals if we are interested in learning how safe the state is. What we
really should be computing is the murders per capita (murders per
100,000 as the unit). To compute this quantity, the powerful vector
arithmetic capabilities of R come in handy.

\subsubsection{Rescaling a vector}\label{rescaling-a-vector}

In R, arithmetic operations on vectors occur \emph{element-wise}. For a
quick example, suppose we have height in inches:

\begin{Shaded}
\begin{Highlighting}[]
\NormalTok{inches }\OtherTok{\textless{}{-}} \FunctionTok{c}\NormalTok{(}\DecValTok{69}\NormalTok{, }\DecValTok{62}\NormalTok{, }\DecValTok{66}\NormalTok{, }\DecValTok{70}\NormalTok{, }\DecValTok{70}\NormalTok{, }\DecValTok{73}\NormalTok{, }\DecValTok{67}\NormalTok{, }\DecValTok{73}\NormalTok{, }\DecValTok{67}\NormalTok{, }\DecValTok{70}\NormalTok{)}
\end{Highlighting}
\end{Shaded}

and want to convert to centimeters. Notice what happens when we multiply
\texttt{inches} by 2.54:

\begin{Shaded}
\begin{Highlighting}[]
\NormalTok{inches }\SpecialCharTok{*} \FloatTok{2.54}
\end{Highlighting}
\end{Shaded}

\begin{verbatim}
##  [1] 175.26 157.48 167.64 177.80 177.80 185.42 170.18 185.42 170.18 177.80
\end{verbatim}

In the line above, we multiplied each element by 2.54. Similarly, if for
each entry we want to compute how many inches taller or shorter than 69
inches, the average height for males, we can subtract it from every
entry like this:

\begin{Shaded}
\begin{Highlighting}[]
\NormalTok{inches }\SpecialCharTok{{-}} \DecValTok{69}
\end{Highlighting}
\end{Shaded}

\begin{verbatim}
##  [1]  0 -7 -3  1  1  4 -2  4 -2  1
\end{verbatim}

\subsubsection{Two vectors}\label{two-vectors}

If we have two vectors of the same length, and we sum them in R, they
will be added entry by entry as follows:

\[
\begin{pmatrix}
a\\
b\\
c\\
d
\end{pmatrix}
+
\begin{pmatrix}
e\\
f\\
g\\
h
\end{pmatrix}
=
\begin{pmatrix}
a +e\\
b + f\\
c + g\\
d + h
\end{pmatrix}
\]

The same holds for other mathematical operations, such as \texttt{-},
\texttt{*} and \texttt{/}.

This implies that to compute the murder rates we can simply type:

\begin{Shaded}
\begin{Highlighting}[]
\NormalTok{murder\_rate }\OtherTok{\textless{}{-}}\NormalTok{ murders}\SpecialCharTok{$}\NormalTok{total }\SpecialCharTok{/}\NormalTok{ murders}\SpecialCharTok{$}\NormalTok{population }\SpecialCharTok{*} \DecValTok{100000}
\end{Highlighting}
\end{Shaded}

Once we do this, we notice that California is no longer near the top of
the list. In fact, we can use what we have learned to order the states
by murder rate:

\begin{Shaded}
\begin{Highlighting}[]
\NormalTok{murders}\SpecialCharTok{$}\NormalTok{abb[}\FunctionTok{order}\NormalTok{(murder\_rate)]}
\end{Highlighting}
\end{Shaded}

\begin{verbatim}
##  [1] "VT" "NH" "HI" "ND" "IA" "ID" "UT" "ME" "WY" "OR" "SD" "MN" "MT" "CO" "WA"
## [16] "WV" "RI" "WI" "NE" "MA" "IN" "KS" "NY" "KY" "AK" "OH" "CT" "NJ" "AL" "IL"
## [31] "OK" "NC" "NV" "VA" "AR" "TX" "NM" "CA" "FL" "TN" "PA" "AZ" "GA" "MS" "MI"
## [46] "DE" "SC" "MD" "MO" "LA" "DC"
\end{verbatim}

\subsection{Indexing}\label{indexing}

R provides a powerful and convenient way of indexing vectors. We can,
for example, subset a vector based on properties of another vector. In
this section, we continue working with our US murders example, which we
can load like this:

\begin{Shaded}
\begin{Highlighting}[]
\FunctionTok{library}\NormalTok{(dslabs)}
\FunctionTok{data}\NormalTok{(}\StringTok{"murders"}\NormalTok{)}
\end{Highlighting}
\end{Shaded}

\subsubsection{Subsetting with logicals}\label{subsetting-with-logicals}

We have now calculated the murder rate using:

\begin{Shaded}
\begin{Highlighting}[]
\NormalTok{murder\_rate }\OtherTok{\textless{}{-}}\NormalTok{ murders}\SpecialCharTok{$}\NormalTok{total }\SpecialCharTok{/}\NormalTok{ murders}\SpecialCharTok{$}\NormalTok{population }\SpecialCharTok{*} \DecValTok{100000} 
\end{Highlighting}
\end{Shaded}

Imagine you are moving from Italy where, according to an ABC news
report, the murder rate is only 0.71 per 100,000. You would prefer to
move to a state with a similar murder rate. Another powerful feature of
R is that we can use logicals to index vectors. If we compare a vector
to a single number, it actually performs the test for each entry. The
following is an example related to the question above:

\begin{Shaded}
\begin{Highlighting}[]
\NormalTok{ind }\OtherTok{\textless{}{-}}\NormalTok{ murder\_rate }\SpecialCharTok{\textless{}} \FloatTok{0.71}
\end{Highlighting}
\end{Shaded}

If we instead want to know if a value is less or equal, we can use:

\begin{Shaded}
\begin{Highlighting}[]
\NormalTok{ind }\OtherTok{\textless{}{-}}\NormalTok{ murder\_rate }\SpecialCharTok{\textless{}=} \FloatTok{0.71}
\end{Highlighting}
\end{Shaded}

Note that we get back a logical vector with \texttt{TRUE} for each entry
smaller than or equal to 0.71. To see which states these are, we can
leverage the fact that vectors can be indexed with logicals.

\begin{Shaded}
\begin{Highlighting}[]
\NormalTok{murders}\SpecialCharTok{$}\NormalTok{state[ind]}
\end{Highlighting}
\end{Shaded}

\begin{verbatim}
## [1] "Hawaii"        "Iowa"          "New Hampshire" "North Dakota" 
## [5] "Vermont"
\end{verbatim}

In order to count how many are TRUE, the function \texttt{sum} returns
the sum of the entries of a vector and logical vectors get
\emph{coerced} to numeric with \texttt{TRUE} coded as 1 and
\texttt{FALSE} as 0. Thus we can count the states using:

\begin{Shaded}
\begin{Highlighting}[]
\FunctionTok{sum}\NormalTok{(ind)}
\end{Highlighting}
\end{Shaded}

\begin{verbatim}
## [1] 5
\end{verbatim}

\subsubsection{Logical operators}\label{logical-operators}

Suppose we like the mountains and we want to move to a safe state in the
western region of the country. We want the murder rate to be at most 1.
In this case, we want two different things to be true. Here we can use
the logical operator \emph{and}, which in R is represented with
\texttt{\&}. This operation results in \texttt{TRUE} only when both
logicals are \texttt{TRUE}. To see this, consider this example:

\begin{Shaded}
\begin{Highlighting}[]
\ConstantTok{TRUE} \SpecialCharTok{\&} \ConstantTok{TRUE}
\end{Highlighting}
\end{Shaded}

\begin{verbatim}
## [1] TRUE
\end{verbatim}

\begin{Shaded}
\begin{Highlighting}[]
\ConstantTok{TRUE} \SpecialCharTok{\&} \ConstantTok{FALSE}
\end{Highlighting}
\end{Shaded}

\begin{verbatim}
## [1] FALSE
\end{verbatim}

\begin{Shaded}
\begin{Highlighting}[]
\ConstantTok{FALSE} \SpecialCharTok{\&} \ConstantTok{FALSE}
\end{Highlighting}
\end{Shaded}

\begin{verbatim}
## [1] FALSE
\end{verbatim}

For our example, we can form two logicals:

\begin{Shaded}
\begin{Highlighting}[]
\NormalTok{west }\OtherTok{\textless{}{-}}\NormalTok{ murders}\SpecialCharTok{$}\NormalTok{region }\SpecialCharTok{==} \StringTok{"West"}
\NormalTok{safe }\OtherTok{\textless{}{-}}\NormalTok{ murder\_rate }\SpecialCharTok{\textless{}=} \DecValTok{1}
\end{Highlighting}
\end{Shaded}

and we can use the \texttt{\&} to get a vector of logicals that tells us
which states satisfy both conditions:

\begin{Shaded}
\begin{Highlighting}[]
\NormalTok{ind }\OtherTok{\textless{}{-}}\NormalTok{ safe }\SpecialCharTok{\&}\NormalTok{ west}
\NormalTok{murders}\SpecialCharTok{$}\NormalTok{state[ind]}
\end{Highlighting}
\end{Shaded}

\begin{verbatim}
## [1] "Hawaii"  "Idaho"   "Oregon"  "Utah"    "Wyoming"
\end{verbatim}

\subsubsection{\texorpdfstring{\texttt{which}}{which}}\label{which}

Suppose we want to look up California's murder rate. For this type of
operation, it is convenient to convert vectors of logicals into indexes
instead of keeping long vectors of logicals. The function \texttt{which}
tells us which entries of a logical vector are TRUE. So we can type:

\begin{Shaded}
\begin{Highlighting}[]
\NormalTok{ind }\OtherTok{\textless{}{-}} \FunctionTok{which}\NormalTok{(murders}\SpecialCharTok{$}\NormalTok{state }\SpecialCharTok{==} \StringTok{"California"}\NormalTok{)}
\NormalTok{murder\_rate[ind]}
\end{Highlighting}
\end{Shaded}

\begin{verbatim}
## [1] 3.374138
\end{verbatim}

\subsubsection{\texorpdfstring{\texttt{match}}{match}}\label{match}

If instead of just one state we want to find out the murder rates for
several states, say New York, Florida, and Texas, we can use the
function \texttt{match}. This function tells us which indexes of a
second vector match each of the entries of a first vector:

\begin{Shaded}
\begin{Highlighting}[]
\NormalTok{ind }\OtherTok{\textless{}{-}} \FunctionTok{match}\NormalTok{(}\FunctionTok{c}\NormalTok{(}\StringTok{"New York"}\NormalTok{, }\StringTok{"Florida"}\NormalTok{, }\StringTok{"Texas"}\NormalTok{), murders}\SpecialCharTok{$}\NormalTok{state)}
\NormalTok{ind}
\end{Highlighting}
\end{Shaded}

\begin{verbatim}
## [1] 33 10 44
\end{verbatim}

Now we can look at the murder rates:

\begin{Shaded}
\begin{Highlighting}[]
\NormalTok{murder\_rate[ind]}
\end{Highlighting}
\end{Shaded}

\begin{verbatim}
## [1] 2.667960 3.398069 3.201360
\end{verbatim}

\subsubsection{\texorpdfstring{\texttt{\%in\%}}{\%in\%}}\label{in}

If rather than an index we want a logical that tells us whether or not
each element of a first vector is in a second, we can use the function
\texttt{\%in\%}. Let's imagine you are not sure if Boston, Dakota, and
Washington are states. You can find out like this:

\begin{Shaded}
\begin{Highlighting}[]
\FunctionTok{c}\NormalTok{(}\StringTok{"Boston"}\NormalTok{, }\StringTok{"Dakota"}\NormalTok{, }\StringTok{"Washington"}\NormalTok{) }\SpecialCharTok{\%in\%}\NormalTok{ murders}\SpecialCharTok{$}\NormalTok{state}
\end{Highlighting}
\end{Shaded}

\begin{verbatim}
## [1] FALSE FALSE  TRUE
\end{verbatim}

Note that we will be using \texttt{\%in\%} often throughout the book.

\textbf{Advanced}: There is a connection between \texttt{match} and
\texttt{\%in\%} through \texttt{which}. To see this, notice that the
following two lines produce the same index (although in different
order):

\begin{Shaded}
\begin{Highlighting}[]
\FunctionTok{match}\NormalTok{(}\FunctionTok{c}\NormalTok{(}\StringTok{"New York"}\NormalTok{, }\StringTok{"Florida"}\NormalTok{, }\StringTok{"Texas"}\NormalTok{), murders}\SpecialCharTok{$}\NormalTok{state)}
\end{Highlighting}
\end{Shaded}

\begin{verbatim}
## [1] 33 10 44
\end{verbatim}

\begin{Shaded}
\begin{Highlighting}[]
\FunctionTok{which}\NormalTok{(murders}\SpecialCharTok{$}\NormalTok{state}\SpecialCharTok{\%in\%}\FunctionTok{c}\NormalTok{(}\StringTok{"New York"}\NormalTok{, }\StringTok{"Florida"}\NormalTok{, }\StringTok{"Texas"}\NormalTok{))}
\end{Highlighting}
\end{Shaded}

\begin{verbatim}
## [1] 10 33 44
\end{verbatim}

\subsection{R programming}\label{r-programming}

We teach R because it greatly facilitates data analysis. By coding in R,
we can efficiently perform exploratory data analysis, build data
analysis pipelines, and prepare data visualization to communicate
results. However, R is not just a data analysis environment but a
programming language. Advanced R programmers can develop complex
packages and even improve R itself, but we do not cover advanced
programming. Nonetheless, in Appendix @ref(appendix-r-programming), we
introduce three key programming concepts: conditional expressions,
for-loops, and functions. These are not just key building blocks for
advanced programming, but are sometimes useful during data analysis.

\pagebreak

\end{document}
